\begin{bt}%[0H9N1-3]
% nguon: Gợi ý từ Ví dụ 4, Bài 10 – SGK Toán 10 KNTT tập 1, trang 63
Trong mặt phẳng toạ độ $Oxy$, cho ba điểm không thẳng hàng $A(1; 3)$, $B(-2; 6)$, $C(5; 1)$.
\begin{enumerate}[a)]
	\item Tìm toạ độ của vectơ $\vec{AB}$.
	\item Tìm toạ độ trung điểm $I$ của đoạn thẳng $AB$.
	\item Tìm toạ độ trọng tâm $G$ của tam giác $ABC$.
\end{enumerate}
\loigiai{
\begin{enumerate}[a)]
	\item Toạ độ của vectơ $\vec{AB}$ là:
	$\vec{AB} = (x_B - x_A; y_B - y_A) = (-2 - 1; 6 - 3) = (-3; 3)$.
	\item Vì $I$ là trung điểm của đoạn thẳng $AB$ nên:
	$x_I = \dfrac{x_A + x_B}{2} = \dfrac{1 + (-2)}{2} = -\dfrac{1}{2}$, $y_I = \dfrac{y_A + y_B}{2} = \dfrac{3 + 6}{2} = \dfrac{9}{2}$.\\
	Vậy $I\left(-\dfrac{1}{2}; \dfrac{9}{2}\right)$.
	\item Vì $G$ là trọng tâm của tam giác $ABC$ nên:
	$x_G = \dfrac{x_A + x_B + x_C}{3} = \dfrac{1 + (-2) + 5}{3} = \dfrac{4}{3}$, $y_G = \dfrac{y_A + y_B + y_C}{3} = \dfrac{3 + 6 + 1}{3} = \dfrac{10}{3}$.\\
	Vậy $G\left(\dfrac{4}{3}; \dfrac{10}{3}\right)$.
\end{enumerate}
}
\end{bt}

\begin{bt}%[0H9H3-2]
% nguon: SBT Toán 10 KNTT tập 2 – Bài 7.55, trang 49, lời giải trang 108
Trong mặt phẳng toạ độ $Oxy$, cho tam giác $ABC$ với $A(1; -1)$, $B(3; 5)$, $C(-2; 4)$.
\begin{enumerate}[a)]
	\item Viết phương trình tham số của đường thẳng $AB$.
	\item Viết phương trình tổng quát của đường cao $AH$ của tam giác $ABC$.
	\item Tính khoảng cách từ điểm $A$ đến đường thẳng $BC$.
	\item Tính $\sin$ của góc giữa hai đường thẳng $AB$ và $AC$.
\end{enumerate}
\loigiai{
\begin{enumerate}[a)]
	\item Ta có $\vec{AB} = (2; 6)$ là một vectơ chỉ phương của đường thẳng $AB$, suy ra $\vec{u} = (1; 3)$ cũng là một vectơ chỉ phương của $AB$.\\
	Đường thẳng $AB$ đi qua điểm $A(1; -1)$ và nhận $\vec{u} = (1; 3)$ làm vectơ chỉ phương nên có phương trình tham số là
	$\begin{cases} x = 1 + t \\ y = -1 + 3t \end{cases}$ ($t \in \mathbb{R}$).
	\item Ta có $\vec{BC} = (-5; -1)$. Vì đường cao $AH$ vuông góc với cạnh $BC$ nên nhận $\vec{n} = (5; 1)$ làm một vectơ pháp tuyến.\\
	Đường thẳng $AH$ đi qua $A(1; -1)$ có phương trình tổng quát là:
	$5(x - 1) + 1(y + 1) = 0 \Leftrightarrow 5x + y - 4 = 0$.
	\item Đường thẳng $BC$ có vectơ chỉ phương $\vec{BC} = (-5; -1)$ nên nhận $\vec{n}' = (1; -5)$ làm một vectơ pháp tuyến. Phương trình tổng quát của đường thẳng $BC$ đi qua $B(3; 5)$ là:
	$1(x - 3) - 5(y - 5) = 0 \Leftrightarrow x - 5y + 22 = 0$.\\
	Khoảng cách từ điểm $A(1; -1)$ đến đường thẳng $BC$ là:
	$d(A, BC) = \dfrac{|1 - 5\cdot(-1) + 22|}{\sqrt{1^2 + (-5)^2}} = \dfrac{28}{\sqrt{26}} = \dfrac{14\sqrt{26}}{13}$.
	\item Gọi $\alpha$ là góc giữa hai đường thẳng $AB$ và $AC$. Ta có $\vec{AB} = (2; 6)$ và $\vec{AC} = (-3; 5)$.\\
	Khi đó:
	$\cos\alpha = |\cos(\vec{AB}, \vec{AC})| = \dfrac{|2\cdot(-3) + 6\cdot 5|}{\sqrt{2^2 + 6^2}\cdot\sqrt{(-3)^2 + 5^2}} = \dfrac{24}{\sqrt{40}\cdot\sqrt{34}} = \dfrac{6}{\sqrt{85}}$.\\
	Suy ra $\sin\alpha = \sqrt{1 - \cos^2\alpha} = \sqrt{1 - \dfrac{36}{85}} = \dfrac{7}{\sqrt{85}} = \dfrac{7\sqrt{85}}{85}$.
\end{enumerate}
}
\end{bt}

\begin{bt}%[0H9H4-3]
% nguon: SBT Toán 10 KNTT tập 2 – Bài 7.57, trang 50, lời giải trang 109
Trong mặt phẳng toạ độ $Oxy$, cho đường tròn $(C)$ có phương trình $x^2 + y^2 - 4x + 6y - 12 = 0$.
\begin{enumerate}[a)]
	\item Tìm toạ độ tâm $I$ và bán kính $R$ của $(C)$.
	\item Chứng minh rằng điểm $M(5; 1)$ thuộc đường tròn $(C)$. Viết phương trình tiếp tuyến $d$ của $(C)$ tại điểm $M$.
\end{enumerate}
\loigiai{
\begin{enumerate}[a)]
	\item Phương trình đường tròn $(C)$ có dạng $x^2 + y^2 - 2ax - 2by + c = 0$ với $a = 2$, $b = -3$, $c = -12$.\\
	Do đó $(C)$ có tâm $I(2; -3)$ và bán kính $R = \sqrt{a^2 + b^2 - c} = \sqrt{2^2 + (-3)^2 - (-12)} = \sqrt{25} = 5$.
	\item Thay toạ độ điểm $M(5; 1)$ vào vế trái phương trình của $(C)$ ta có:
	$5^2 + 1^2 - 4\cdot 5 + 6\cdot 1 - 12 = 25 + 1 - 20 + 6 - 12 = 0$.\\
	Vậy điểm $M(5; 1)$ thuộc đường tròn $(C)$.\\
	Tiếp tuyến $d$ của đường tròn $(C)$ tại điểm $M$ đi qua $M(5; 1)$ và vuông góc với bán kính $IM$ nên nhận vectơ $\vec{IM} = (5 - 2; 1 - (-3)) = (3; 4)$ làm một vectơ pháp tuyến.\\
	Phương trình tiếp tuyến $d$ là:
	$3(x - 5) + 4(y - 1) = 0 \Leftrightarrow 3x + 4y - 19 = 0$.
\end{enumerate}
}
\end{bt}

\begin{bt}%[0H9V5-0]
% nguon: SBT Toán 10 KNTT tập 2 – Bài 7.35, trang 46, lời giải trang 105–106
Một người kĩ sư thiết kế một đường hầm một chiều có mặt cắt là một nửa hình elip, chiều rộng của hầm là $12\text{ m}$, khoảng cách từ điểm cao nhất của elip so với mặt đường là $3\text{ m}$. Người kĩ sư này muốn đưa ra cảnh báo cho các xe có thể đi qua hầm. Biết rằng những loại xe tải có chiều cao $2{,}8\text{ m}$ thì có chiều rộng không quá $3\text{ m}$. Hỏi chiếc xe tải có chiều cao $2{,}8\text{ m}$ có thể đi qua hầm an toàn được không?
\loigiai{
Chọn hệ trục toạ độ $Oxy$ sao cho trục $Ox$ nằm dọc theo mặt đường, gốc $O$ là trung điểm đoạn thẳng nối hai chân vòm hầm trên mặt đường, trục $Oy$ hướng thẳng đứng lên trên (trùng với trục đối xứng của mặt cắt).\\
Khi đó nửa hình elip nằm phía trên trục $Ox$ thuộc elip $(E)$ có phương trình chính tắc dạng $\dfrac{x^2}{a^2} + \dfrac{y^2}{b^2} = 1$ ($a > b > 0$).\\
Vì chiều rộng của hầm là $12\text{ m}$ và khoảng cách từ điểm cao nhất của elip so với mặt đường là $3\text{ m}$ nên elip $(E)$ có độ dài trục lớn là $2a = 12 \Rightarrow a = 6$ và nửa trục bé là $b = 3$.\\
Do đó phương trình chính tắc của elip $(E)$ là: $\dfrac{x^2}{36} + \dfrac{y^2}{9} = 1$.\\
Để đi qua hầm an toàn nhất, xe tải cần chạy vào chính giữa hầm. Với xe tải có chiều rộng không quá $3\text{ m}$, điểm cao nhất ở mép ngoài của thân xe ứng với hoành độ $x = 1{,}5$.\\
Thay $x = 1{,}5$ vào phương trình elip $(E)$ với tung độ $y > 0$, ta tính được độ cao tối đa của vòm hầm tại vị trí này là:
$y_M = 3\sqrt{1 - \dfrac{1{,}5^2}{36}} = 3\sqrt{1 - \dfrac{1}{16}} = \dfrac{3\sqrt{15}}{4} \approx 2{,}905\text{ (m)}$.\\
Vì $2{,}905\text{ m} > 2{,}8\text{ m}$ nên chiếc xe tải có chiều cao $2{,}8\text{ m}$ có thể đi qua hầm an toàn được khi chạy ở chính giữa hầm.
}
\end{bt}

\begin{bt}%[0H9C3-6]
% nguon: Trích và mở rộng ý b từ Bài 7.17 – SBT Toán 10 KNTT tập 2, trang 38–39
Trong mặt phẳng toạ độ $Oxy$, cho hai điểm $A(-3; 0)$, $B(1; -2)$ và đường thẳng $d: x + y - 1 = 0$.
\begin{enumerate}[a)]
	\item Chứng minh rằng hai điểm $A$ và $B$ nằm cùng phía so với đường thẳng $d$.
	\item Điểm $M$ thay đổi trên đường thẳng $d$. Tìm toạ độ điểm $M$ để chu vi tam giác $ABM$ đạt giá trị nhỏ nhất và tính giá trị nhỏ nhất đó.
\end{enumerate}
\loigiai{
\begin{enumerate}[a)]
	\item Đặt $f(x; y) = x + y - 1$. Thay toạ độ các điểm $A(-3; 0)$ và $B(1; -2)$ vào $f(x; y)$ ta có:
	$f(-3; 0) = -3 + 0 - 1 = -4$, $f(1; -2) = 1 + (-2) - 1 = -2$.\\
	Vì $f(-3; 0)\cdot f(1; -2) = (-4)\cdot(-2) = 8 > 0$ nên hai điểm $A$ và $B$ nằm cùng phía đối với đường thẳng $d$.
	\item Chu vi tam giác $ABM$ là $P_{\triangle ABM} = AB + MA + MB$.\\
	Độ dài đoạn thẳng $AB$ là:
	$AB = \sqrt{(1 - (-3))^2 + (-2 - 0)^2} = \sqrt{16 + 4} = 2\sqrt{5}$ (không đổi).\\
	Do đó chu vi tam giác $ABM$ nhỏ nhất khi và chỉ khi tổng $MA + MB$ đạt giá trị nhỏ nhất.\\
	Gọi $A'$ là điểm đối xứng với điểm $A(-3; 0)$ qua đường thẳng $d$.\\
	Đường thẳng $\Delta$ đi qua $A(-3; 0)$ và vuông góc với $d$ có vectơ chỉ phương của $d$ là $\vec{u}_d = (1; -1)$ làm vectơ pháp tuyến nên có phương trình:
	$1(x + 3) - 1(y - 0) = 0 \Leftrightarrow x - y + 3 = 0$.\\
	Hình chiếu vuông góc $H$ của $A$ lên $d$ là nghiệm của hệ phương trình:
	$\begin{cases} x + y - 1 = 0 \\ x - y + 3 = 0 \end{cases} \Leftrightarrow \begin{cases} x = -1 \\ y = 2 \end{cases} \Rightarrow H(-1; 2)$.\\
	Vì $H$ là trung điểm của đoạn thẳng $AA'$ nên:
	$\begin{cases} x_{A'} = 2x_H - x_A = 2(-1) - (-3) = 1 \\ y_{A'} = 2y_H - y_A = 2\cdot 2 - 0 = 4 \end{cases} \Rightarrow A'(1; 4)$.\\
	Do $A'$ đối xứng với $A$ qua đường thẳng $d$ và $M \in d$ nên $MA = MA'$.\\
	Áp dụng bất đẳng thức tam giác ta có:
	$MA + MB = MA' + MB \ge A'B$.\\
	Vì $A$ và $B$ nằm cùng một phía so với đường thẳng $d$ nên $A'$ và $B$ nằm về hai phía khác nhau so với đường thẳng $d$. Do đó đoạn thẳng $A'B$ cắt đường thẳng $d$.\\
	Dấu bằng xảy ra khi và chỉ khi $M$ là giao điểm của đoạn thẳng $A'B$ với đường thẳng $d$.\\
	Hai điểm $A'(1; 4)$ và $B(1; -2)$ có cùng hoành độ bằng $1$ nên đường thẳng $A'B$ có phương trình $x = 1$.\\
	Toạ độ điểm $M$ là nghiệm của hệ phương trình:
	$\begin{cases} x = 1 \\ x + y - 1 = 0 \end{cases} \Leftrightarrow \begin{cases} x = 1 \\ y = 0 \end{cases} \Rightarrow M(1; 0)$.\\
	Khi đó: $A'B = \sqrt{(1 - 1)^2 + (-2 - 4)^2} = \sqrt{0 + 36} = 6$.\\
	Vậy chu vi tam giác $ABM$ đạt giá trị nhỏ nhất bằng $6 + 2\sqrt{5}$ khi $M(1; 0)$.
\end{enumerate}
}
\end{bt}
